\documentclass[11pt,a4paper]{article}

% =========================================================
% Packages
% =========================================================
\usepackage[margin=1in]{geometry}
\usepackage{times}
\usepackage{microtype}

\usepackage{amsmath,amssymb,amsfonts,amsthm}
\usepackage{booktabs}
\usepackage{multirow}
\usepackage{array}
\usepackage{adjustbox}
\usepackage{enumitem}

\usepackage{algorithm}
\usepackage{algorithmic}

\usepackage{graphicx}
\usepackage{svg}
\usepackage{subcaption}
\usepackage{float}
\usepackage{placeins}

\usepackage{xcolor}
\usepackage{url}
\usepackage[colorlinks=true,linkcolor=blue,citecolor=blue,urlcolor=blue]{hyperref}

% =========================================================
% Theorem Environments
% =========================================================
\newtheorem{definition}{Definition}
\newtheorem{theorem}{Theorem}
\newtheorem{lemma}{Lemma}
\newtheorem{property}{Property}
\newtheorem{corollary}{Corollary}
\newtheorem{remark}{Remark}

% =========================================================
% Custom Commands
% =========================================================
\newcommand{\D}{\mathcal{D}}
\newcommand{\I}{\mathcal{I}}
\newcommand{\W}{\mathcal{W}}
\newcommand{\HIEP}{\mathcal{H}}
\newcommand{\Visited}{\mathcal{V}}
\newcommand{\Output}{\mathcal{O}}
\newcommand{\Signal}{\mathcal{S}}
\newcommand{\Junk}{\mathcal{J}}

\newcommand{\supp}{\operatorname{sup}}
\newcommand{\df}{\operatorname{df}}
\newcommand{\prob}{\operatorname{Pr}}
\newcommand{\spanop}{\operatorname{span}}
\newcommand{\uniq}{\operatorname{uniq}}
\newcommand{\pwt}{\operatorname{pwt}}
\newcommand{\tail}{\operatorname{tail}}
\newcommand{\TIUB}{\operatorname{TIUB}}
\newcommand{\ITU}{\operatorname{ITU}}
\newcommand{\IWRU}{\operatorname{IWRU}}
\newcommand{\JOIN}{\operatorname{JOIN}}
\newcommand{\RSS}{\operatorname{RSS}}
\newcommand{\NFE}{\operatorname{NFE}}
\newcommand{\IDO}{\operatorname{IDO}}

% =========================================================
% Title
% =========================================================
\title{\textbf{HIEP-Miner: High-Information Entropy Pattern Mining for Embedded Text Streams}}

\author{
Quan Van\\
Independent Researcher\\
\texttt{vanhaminhquan2406@gmail.com}
}

\date{}

\begin{document}

\maketitle

% =========================================================
% Abstract
% =========================================================
\begin{abstract}
Classical pattern mining is primarily support-driven. While support enables powerful downward-closure pruning, it is often semantically misaligned with text mining because natural language follows a heavy-tailed Zipfian distribution: high-support tokens are frequently stopwords, boilerplate terms, or workflow markers rather than meaningful concepts. This paper introduces \emph{High-Information Entropy Pattern Mining} (HIEPM), a new mining paradigm for embedded text-stream engines that mines rare but recurring, highly informative token combinations from integerized text streams. HIEPM uses Shannon self-information as a token-level utility signal and defines a support-weighted information utility for itemset and sequence patterns extracted from window-bounded text transactions.

The proposed model is designed for a zero-dependency C99-style engine using memory-mapped input, a flat hash tokenizer, and integer-only mining structures. We show that direct joint-surprisal objectives are neither monotone nor anti-monotone, making classical Apriori, FP-Growth, and PrefixSpan pruning insufficient. We therefore introduce a decomposable entropy-derived utility objective and derive two safe pruning bounds: Transaction Information Upper Bound (TIUB) and Information-Weighted Remaining Utility (IWRU). The IWRU theorem provides a monotone branch-and-bound envelope for exact high-information pattern mining. We then propose HIEP-Miner, a cache-friendly vertical utility-list algorithm using flat arrays, integer postings, information-weighted suffix arrays, and depth-first arena reuse. The paper provides formal definitions, theoretical guarantees, complexity analysis, correctness proofs, and an executed experimental evaluation using Retail, Accidents, Chess, Databricks Dolly text, and planted Zipfian text streams.
\end{abstract}

\noindent\textbf{Keywords:}
information theory; self-information; entropy; text mining; high-utility itemset mining; sequential pattern mining; mmap; C99; pattern mining; Zipfian text.

% =========================================================
\section{Introduction}
% =========================================================

Pattern mining traditionally relies on support. Apriori uses support anti-monotonicity to prune candidate itemsets level by level \cite{AgrawalSrikant1994}. FP-Growth reduces candidate generation using a compressed FP-tree representation \cite{HanPeiYin2000}. Eclat uses vertical transaction-id sets for efficient support computation \cite{Zaki2000}. PrefixSpan applies prefix-projection to sequential pattern mining \cite{Pei2001PrefixSpan}. These algorithms remain foundational because support has a clean downward-closure property: if a pattern is infrequent, then all of its super-patterns are infrequent.

However, support is often semantically weak for text. Natural language follows a Zipfian distribution, where a small number of tokens occur extremely frequently and a long tail of tokens occurs rarely \cite{Zipf1949,Piantadosi2014Zipf}. In text streams, high-support patterns are often dominated by function words, workflow labels, conversational markers, or boilerplate terms such as ``please'', ``system'', ``user'', and ``assistant''. These tokens can dominate frequent pattern outputs while contributing little semantic specificity.

Information retrieval addressed this problem at the term level through inverse document frequency. The idf weight of a term increases as the document frequency of the term decreases, thereby downweighting ubiquitous terms \cite{ManningRaghavanSchutze2008IR}. From an information-theoretic perspective, this corresponds to assigning higher self-information to rarer events. Shannon's self-information of an event \(x\) is:
\[
I(x)=-\log p(x),
\]
where \(p(x)\) is the probability of observing \(x\) \cite{Shannon1948}. This principle suggests a different mining objective: instead of mining merely frequent token combinations, mine recurring token combinations with high information content.

Several research lines are related to this idea. Information-theoretic pattern mining has studied maximum-entropy significance, subjective interestingness, and exploratory data mining based on information gain \cite{Tatti2008MaximumEntropy,DeBie2011SubjectiveInterestingness,Kontonasios2012SubjectiveInterestingness}. Minimum Description Length (MDL) methods such as Krimp mine compact pattern sets that compress data effectively \cite{Vreeken2011Krimp,Galbrun2020MDL}. High-utility itemset mining (HUIM) also addresses non-support objectives by introducing upper bounds such as transaction-weighted utilization, utility lists, remaining utility, and subtree utility \cite{Liu2005TwoPhase,Liu2012HUIMiner,FournierViger2014FHM,Zida2017EFIM}. Nevertheless, these approaches do not directly provide an embedded, exact, high-throughput miner for entropy-derived token patterns over integerized text streams.

This paper proposes \emph{High-Information Entropy Pattern Mining} (HIEPM). The term ``entropy'' is used to denote a Shannon-inspired information-theoretic mining paradigm. Strictly, entropy is an expectation over a distribution, while the information content of a realized token is self-information or surprisal. Accordingly, HIEPM mines high self-information patterns and reports entropy-derived utility values.

The proposed algorithm, HIEP-Miner, is designed for an embedded text mining engine called \texttt{dm.exe}. The assumed pipeline is:
\[
\text{raw text}
\rightarrow
\text{memory-mapped scan}
\rightarrow
\text{flat hash tokenization}
\rightarrow
\text{integer token stream}
\rightarrow
\text{window-bounded transactions}
\rightarrow
\text{HIEP-Miner}.
\]

The contributions of this paper are:

\begin{itemize}
    \item We define the High-Information Entropy Pattern Mining problem for integerized text streams.
    \item We formalize self-information token weights, window-bounded text transactions, itemset and sequence modes, compactness-aware utility, and high-information pattern outputs.
    \item We show why direct joint-surprisal objectives break monotone pruning.
    \item We introduce Transaction Information Upper Bound (TIUB), a coarse entropy-derived transaction upper bound.
    \item We prove the Information-Weighted Remaining Utility (IWRU) theorem, which gives a safe and monotone pruning envelope for exact mining.
    \item We propose HIEP-Miner, a flat-array vertical utility-list miner suitable for zero-dependency C99 implementation.
    \item We provide correctness proofs, output-sensitive complexity analysis, and a reproducible experimental validation using real and synthetic datasets.
\end{itemize}

% =========================================================
\section{Related Work}
% =========================================================

\subsection{Support-Based Pattern Mining}

Apriori established frequent itemset mining based on support anti-monotonicity \cite{AgrawalSrikant1994}. FP-Growth introduced the FP-tree and avoided explicit candidate generation \cite{HanPeiYin2000}. Eclat popularized vertical transaction-id set intersections \cite{Zaki2000}. PrefixSpan extended the prefix-growth philosophy to sequential pattern mining \cite{Pei2001PrefixSpan}. These algorithms are effective when support is aligned with the user's notion of interestingness. In text, however, high support often indicates genericity rather than informativeness.

\subsection{Information Theory in Text Mining}

Shannon's information theory defines entropy as expected information and self-information as the surprise of a realized event \cite{Shannon1948}. In information retrieval, inverse document frequency downweights common terms and upweights rare terms \cite{ManningRaghavanSchutze2008IR}. Church and Gale related inverse document frequency to deviations from simple random models of word occurrence \cite{ChurchGale1995}. Keyword extraction research also combines phrase frequency and informativeness to identify meaningful phrases \cite{TomokiyoHurst2003}.

HIEPM differs from classical information retrieval because it does not score individual terms only. It mines recurring token combinations whose aggregate information utility is high.

\subsection{Information-Theoretic and MDL Pattern Mining}

Information-theoretic pattern mining has studied itemset significance under maximum-entropy models \cite{Tatti2008MaximumEntropy}. De Bie formulated exploratory data mining as information exchange relative to background knowledge \cite{DeBie2011SubjectiveInterestingness}. Kontonasios and De Bie further studied subjective interestingness of pattern sets \cite{Kontonasios2012SubjectiveInterestingness}. MDL-based approaches such as Krimp select pattern sets that compress data well \cite{Vreeken2011Krimp}, and surveys have summarized the broader role of MDL in pattern mining \cite{Galbrun2020MDL}.

These methods are conceptually close to HIEPM, but they do not target an embedded text-stream setting with integer token streams, memory-mapped ingestion, and low-level exact mining.

\subsection{High-Utility Itemset Mining}

HUIM discovers itemsets whose utility exceeds a user-specified threshold. Since exact utility is not anti-monotone, Two-Phase introduced transaction-weighted utilization (TWU) as an upper bound \cite{Liu2005TwoPhase}. HUI-Miner introduced utility lists to avoid candidate generation \cite{Liu2012HUIMiner}. FHM tightened pruning using estimated utility co-occurrence \cite{FournierViger2014FHM}. EFIM introduced stronger upper bounds, transaction merging, and projection \cite{Zida2017EFIM}. Utility-oriented pattern mining surveys describe this development in detail \cite{Gan2021UtilitySurvey}.

HIEPM borrows the HUIM idea that non-support objectives can be mined exactly if safe upper bounds are designed. However, HIEPM uses entropy-derived token weights rather than economic utility values.

\subsection{Analytical Gap}

Existing methods leave a gap. Support-based miners over-report high-frequency linguistic noise. HUIM methods require static item utility definitions and are not designed around Shannon self-information, window compactness, and token stream execution. Information-theoretic pattern mining often focuses on significance scoring, background models, or compression, rather than exact embedded mining over integerized text streams.

HIEPM fills this gap by combining self-information weighting, decomposable utility, safe upper bounds, and low-level vertical mining structures.

% =========================================================
\section{Problem Formulation}
% =========================================================

\subsection{Token Stream and Windowization}

Let a tokenizer emit an integer token stream:
\[
Z=\langle z_1,z_2,\dots,z_N\rangle,
\qquad
z_j\in\I,
\]
where \(\I=\{1,2,\dots,m\}\) is the token universe produced by the flat hash tokenizer.

Let \(\mathcal{W}_{L,\Delta}\) be a windowization operator with window length \(L\) and stride \(\Delta\). It produces \(n\) windows:
\[
\D=\{T_1,T_2,\dots,T_n\}.
\]

For window \(q\), let:
\[
a_q=1+(q-1)\Delta,
\qquad
b_q=\min(a_q+L-1,N).
\]

In itemset mode, the transaction is:
\[
T_q=\uniq(\langle z_{a_q},z_{a_q+1},\dots,z_{b_q}\rangle)\subseteq\I.
\]

In sequence mode, the transaction is the ordered sequence:
\[
S_q=\langle z_{a_q},z_{a_q+1},\dots,z_{b_q}\rangle.
\]

We write:
\[
X\preceq T
\]
to mean either \(X\subseteq T\) in itemset mode or \(X\) occurs as an ordered subsequence in sequence mode.

\begin{definition}[Supporting Window Set]
The supporting window set of a pattern \(X\) is:
\[
\Gamma(X)=\{T\in\D\mid X\preceq T\}.
\]
\end{definition}

\begin{definition}[Window Support]
The support of \(X\) is:
\[
\supp(X)=|\Gamma(X)|.
\]
\end{definition}

\subsection{Token Probability and Self-Information}

For a singleton token \(x\in\I\), define:
\[
\df(x)=\supp(\{x\}).
\]

With smoothing parameter \(\alpha>0\), define the empirical window probability:
\[
p(x)=\frac{\df(x)+\alpha}{n+\alpha|\I|}.
\]

\begin{definition}[Self-Information Token Weight]
The self-information weight of token \(x\) is:
\[
w(x)=I(x)=-\log_2 p(x).
\]
\end{definition}

Thus, frequent tokens obtain small weights, while rare tokens obtain large weights.

\begin{remark}
The weight \(w(x)\) is corpus-relative. If the corpus distribution changes, the self-information weights should be recomputed or incrementally maintained.
\end{remark}

\subsection{Joint Surprisal and Non-Monotonicity}

A natural but difficult pattern score is empirical joint surprisal:
\[
J(X)=
-\log_2\left(
\frac{\supp(X)+\alpha}{n+\alpha}
\right).
\]

A support-weighted joint-surprisal score is:
\[
S_J(X)=\supp(X)J(X).
\]

\begin{theorem}[Joint Surprisal Non-Monotonicity]
The function \(S_J(X)\) is neither monotone nor anti-monotone over the pattern lattice.
\end{theorem}

\begin{proof}
Let \(X\subset Y\). Since support is anti-monotone, \(\supp(Y)\le\supp(X)\). Therefore \(J(Y)\ge J(X)\), because rarer patterns have larger surprisal. However, \(S_J(Y)=\supp(Y)J(Y)\) multiplies a decreasing factor by an increasing factor. Depending on the rate of support decrease and surprisal increase, \(S_J(Y)\) can be larger or smaller than \(S_J(X)\). Hence the score is neither monotone nor anti-monotone.
\end{proof}

This invalidates direct Apriori-style pruning based on the final score.

\subsection{Entropy-Derived Decomposable Utility}

To obtain an exact mineable objective, HIEPM defines a decomposable semantic information utility.

\begin{definition}[Pattern Information Weight]
For a pattern \(X\), define:
\[
\pwt(X)=\sum_{x\in X}w(x).
\]
\end{definition}

For itemset mode:
\[
\kappa(X,T)=1.
\]

For sequence mode, let \(X=\langle x_1,x_2,\dots,x_k\rangle\). Define:
\[
\spanop_T(X)
=
\min_{\pi_1<\pi_2<\cdots<\pi_k,\;T[\pi_j]=x_j}
(\pi_k-\pi_1+1).
\]

Then define compactness:
\[
\kappa(X,T)=
\exp\left(
-\gamma(\spanop_T(X)-|X|)
\right),
\qquad
\gamma\ge0.
\]

\begin{definition}[Transaction-Level Information Utility]
The information utility of \(X\) in transaction \(T\) is:
\[
u(X,T)=
\mathbf{1}[X\preceq T]\,
\kappa(X,T)
\sum_{x\in X}w(x).
\]
\end{definition}

\begin{definition}[Database-Level Information Utility]
The global information utility of \(X\) is:
\[
U(X)=\sum_{T\in\D}u(X,T).
\]
\end{definition}

In itemset mode:
\[
U(X)=\supp(X)\sum_{x\in X}w(x).
\]

\subsection{High-Information Entropy Pattern Mining}

Let \(\theta>0\) be a minimum information-utility threshold and \(\sigma_0\ge1\) be a minimum persistence threshold.

\begin{definition}[High-Information Entropy Pattern]
A non-empty pattern \(X\) is a High-Information Entropy Pattern if:
\[
U(X)\ge\theta
\]
and
\[
\supp(X)\ge\sigma_0.
\]
\end{definition}

\begin{definition}[HIEPM Problem]
The High-Information Entropy Pattern Mining problem is to discover:
\[
\HIEP_{\theta,\sigma_0}
=
\{X\neq\emptyset\mid U(X)\ge\theta,\ \supp(X)\ge\sigma_0\}.
\]
\end{definition}

The persistence threshold \(\sigma_0\) prevents the output from being dominated by one-off rare artifacts, spelling errors, or accidental token combinations.

% =========================================================
\section{Upper Bounds and Pruning Theory}
% =========================================================

\subsection{Support Anti-Monotonicity}

\begin{theorem}[Support Anti-Monotonicity]
If \(X\preceq Y\), then:
\[
\supp(Y)\le\supp(X).
\]
\end{theorem}

\begin{proof}
Every transaction containing or matching \(Y\) also contains or matches \(X\). Thus:
\[
\Gamma(Y)\subseteq\Gamma(X),
\]
and the result follows by taking cardinalities.
\end{proof}

\begin{property}[Safe Support Pruning]
If:
\[
\supp(X)<\sigma_0,
\]
then no extension of \(X\) can be a High-Information Entropy Pattern.
\end{property}

\subsection{Transaction Information Upper Bound}

\begin{definition}[Transaction Information Utility]
The total information weight of a transaction \(T\) is:
\[
ITU(T)=\sum_{x\in T}w(x).
\]
\end{definition}

\begin{definition}[Transaction Information Upper Bound]
For a pattern \(X\), define:
\[
TIUB(X)=
\sum_{T\in\Gamma(X)}
ITU(T).
\]
\end{definition}

\begin{theorem}[TIUB Upper Bound]
For every extension \(Y\succeq X\):
\[
U(Y)\le TIUB(X).
\]
\end{theorem}

\begin{proof}
If \(Y\) occurs in \(T\), then \(X\) also occurs in \(T\). Moreover:
\[
u(Y,T)
=
\kappa(Y,T)\sum_{y\in Y}w(y)
\le
\sum_{y\in Y}w(y)
\le
ITU(T).
\]
Since \(\Gamma(Y)\subseteq\Gamma(X)\), summing over \(\Gamma(X)\) gives:
\[
U(Y)\le TIUB(X).
\]
\end{proof}

\begin{property}[Safe TIUB Pruning]
If:
\[
TIUB(X)<\theta,
\]
then no extension of \(X\) can satisfy the HIEPM utility threshold.
\end{property}

\subsection{Information-Weighted Remaining Utility}

Fix a global token order \(\prec\). For a prefix \(X\), let \(E_X\) be the suffix extension set under this order. For a transaction \(T\in\Gamma(X)\), define:
\[
\tail_T(X)=E_X\cap T
\]
in itemset mode. In sequence mode, \(\tail_T(X)\) consists of extension-eligible tokens that occur after the last matched position of \(X\).

\begin{definition}[Remaining Information Utility]
For \(T\in\Gamma(X)\), define:
\[
ru(X,T)=\sum_{y\in\tail_T(X)}w(y).
\]
\end{definition}

\begin{definition}[Information-Weighted Remaining Utility]
\[
IWRU(X)=
\sum_{T\in\Gamma(X)}
\left(
\pwt(X)+ru(X,T)
\right).
\]
\end{definition}

\begin{theorem}[Information-Weighted Remaining Utility Bound]
For every extension \(Y\succeq X\):
\[
U(Y)\le IWRU(X).
\]
\end{theorem}

\begin{proof}
Let \(Y=X\cup Z\), where \(Z\subseteq E_X\). For every \(T\in\Gamma(Y)\), the additional items in \(Z\) must be contained in \(\tail_T(X)\). Since \(\kappa(Y,T)\le1\),
\[
u(Y,T)
=
\kappa(Y,T)
\sum_{y\in Y}w(y)
\le
\sum_{y\in Y}w(y).
\]
Also:
\[
\sum_{y\in Y}w(y)
=
\pwt(X)+\sum_{z\in Z}w(z)
\le
\pwt(X)+ru(X,T).
\]
Therefore:
\[
u(Y,T)\le \pwt(X)+ru(X,T).
\]
Since \(\Gamma(Y)\subseteq\Gamma(X)\),
\[
U(Y)
=
\sum_{T\in\Gamma(Y)}u(Y,T)
\le
\sum_{T\in\Gamma(X)}
\left(
\pwt(X)+ru(X,T)
\right)
=
IWRU(X).
\]
\end{proof}

\begin{corollary}[Safe IWRU Pruning]
If:
\[
IWRU(X)<\theta,
\]
then no extension of \(X\) can be a High-Information Entropy Pattern.
\end{corollary}

\begin{theorem}[Monotone Branch Envelope]
For any extension \(Y\succeq X\):
\[
IWRU(Y)\le IWRU(X).
\]
\end{theorem}

\begin{proof}
Along a depth-first extension branch, the supporting transaction set can only shrink:
\[
\Gamma(Y)\subseteq\Gamma(X).
\]
Moreover, the admissible extension tail inside every surviving transaction also shrinks under the fixed global order. Thus, the sum of prefix weight plus remaining tail weight over the surviving transactions cannot increase. Therefore:
\[
IWRU(Y)\le IWRU(X).
\]
\end{proof}

This theorem provides the downward-closure analogue needed for entropy-derived utility mining.

% =========================================================
\section{HIEP-Miner Architecture}
% =========================================================

\subsection{System Model}

HIEP-Miner is designed for an embedded C99-style engine. The assumed system architecture is:

\begin{enumerate}[label=(\arabic*)]
    \item memory-map the raw text file;
    \item tokenize the stream into integer IDs using a flat hash tokenizer;
    \item windowize the ID stream into itemset or sequence transactions;
    \item build vertical postings and information weights;
    \item mine high-information patterns with HIEP-Miner.
\end{enumerate}

The design avoids pointer-heavy trees and dynamic per-node allocation. Instead, it uses contiguous arrays, offsets, and recyclable arenas.

\subsection{Core Data Structures}

\begin{table}[H]
\centering
\caption{Main flat-array structures used by HIEP-Miner.}
\label{tab:structures}
\small
\begin{adjustbox}{max width=\textwidth}
\begin{tabular}{lll}
\toprule
Structure & Logical Fields & Purpose \\
\midrule
Item metadata & \(df, w, TIUB_1, post\_off, post\_len\) & token statistics and singleton bounds \\
Transaction store & \(txn\_off, txn\_items\) & compact window-bounded transactions \\
Vertical postings & \(occ\_tid, occ\_pos\) & singleton occurrence lists \\
Utility-list arena & \(ul\_tid, ul\_last, ul\_eu, ul\_ru\) & depth-first mining workspace \\
Output buffer & pattern IDs and scores & emitted high-information patterns \\
\bottomrule
\end{tabular}
\end{adjustbox}
\end{table}

The layout is structure-of-arrays rather than array-of-structures. This improves sequential memory access and allows each mining pass to read only the fields it needs.

\subsection{Utility-List Entry}

\begin{definition}[HIEP Utility-List Entry]
For a pattern \(X\) and supporting transaction \(T\), an entry is:
\[
e=\langle tid,last,eu,ru\rangle,
\]
where:
\begin{itemize}
    \item \(tid\) is the transaction identifier;
    \item \(last\) is the last matched token position in sequence mode;
    \item \(eu=u(X,T)\) is the exact information utility contribution;
    \item \(ru=ru(X,T)\) is the remaining information utility.
\end{itemize}
\end{definition}

For itemset mode, \(last\) can be omitted or interpreted as a sorted-transaction offset.

For a utility list \(UL(X)\):
\[
\supp(X)=|UL(X)|,
\]
\[
U(X)=\sum_{e\in UL(X)}e.eu,
\]
and:
\[
IWRU(X)=\sum_{e\in UL(X)}\left(\pwt(X)+e.ru\right).
\]

\subsection{Item Ordering}

HIEP-Miner uses the order:
\[
x\prec y
\]
if:
\[
TIUB(\{x\})<TIUB(\{y\}),
\]
with ties broken by:
\[
w(x)>w(y).
\]

This order tends to place low-opportunity tokens earlier while preserving high-information rare tokens as strong discriminators.

% =========================================================
\section{HIEP-Miner Algorithm}
% =========================================================

\begin{algorithm}[H]
\caption{HIEP-Miner}
\label{alg:hiepminer}
\begin{algorithmic}[1]
\REQUIRE Integer token stream \(Z\), window length \(L\), stride \(\Delta\), threshold \(\theta\), persistence floor \(\sigma_0\)
\ENSURE High-Information Entropy Pattern set \(\HIEP_{\theta,\sigma_0}\)
\STATE Construct window-bounded transaction database \(\D=\mathcal{W}_{L,\Delta}(Z)\)
\STATE Compute \(df(x)\), \(p(x)\), and \(w(x)\) for all tokens \(x\in\I\)
\STATE Compute \(ITU(T)\) for all transactions \(T\in\D\)
\STATE Build singleton vertical postings and singleton utility lists
\STATE Compute \(TIUB(\{x\})\) for all singleton tokens
\STATE Remove tokens \(x\) with \(df(x)<\sigma_0\) or \(TIUB(\{x\})<\theta\)
\STATE Sort remaining tokens by ascending \(TIUB\), then descending \(w\)
\STATE \(\Output\leftarrow\emptyset\)
\FOR{each remaining token \(x\) in order}
    \STATE \(X\leftarrow\{x\}\)
    \STATE Compute exact \(U(X)\) from \(UL(X)\)
    \IF{\(U(X)\ge\theta\)}
        \STATE Add \(X\) to \(\Output\)
    \ENDIF
    \IF{\(IWRU(X)\ge\theta\)}
        \STATE \textsc{DFS-HIEP}\((X,UL(X),E_X,\Output)\)
    \ENDIF
\ENDFOR
\RETURN \(\Output\)
\end{algorithmic}
\end{algorithm}

\begin{algorithm}[H]
\caption{\textsc{DFS-HIEP}}
\label{alg:dfs}
\begin{algorithmic}[1]
\REQUIRE Prefix \(X\), utility list \(UL(X)\), suffix item set \(E_X\), output set \(\Output\)
\FOR{each token \(a\in E_X\)}
    \STATE \(Y\leftarrow X\cup\{a\}\)
    \STATE \(UL(Y)\leftarrow\textsc{Join-HIEP}(UL(X),UL(a))\)
    \IF{\(|UL(Y)|<\sigma_0\)}
        \STATE \textbf{continue}
    \ENDIF
    \STATE Compute exact \(U(Y)=\sum_{e\in UL(Y)}e.eu\)
    \IF{\(U(Y)\ge\theta\)}
        \STATE Add \(Y\) to \(\Output\)
    \ENDIF
    \STATE Compute \(IWRU(Y)=\sum_{e\in UL(Y)}(\pwt(Y)+e.ru)\)
    \IF{\(IWRU(Y)\ge\theta\)}
        \STATE \(E_Y\leftarrow\) items after \(a\) in \(E_X\)
        \STATE \textsc{DFS-HIEP}\((Y,UL(Y),E_Y,\Output)\)
    \ENDIF
\ENDFOR
\end{algorithmic}
\end{algorithm}

\begin{algorithm}[H]
\caption{\textsc{Join-HIEP}}
\label{alg:join}
\begin{algorithmic}[1]
\REQUIRE Utility list \(UL(X)\), singleton or derived utility list \(UL(a)\)
\ENSURE Utility list \(UL(X\cup\{a\})\)
\STATE \(UL(Y)\leftarrow\emptyset\)
\STATE Merge \(UL(X)\) and \(UL(a)\) by transaction identifier
\FOR{each common transaction \(T\)}
    \IF{sequence mode and \(a\) does not occur after the last position of \(X\)}
        \STATE \textbf{continue}
    \ENDIF
    \STATE Compute compactness \(\kappa(Y,T)\)
    \STATE \(eu\leftarrow \kappa(Y,T)\pwt(Y)\)
    \STATE \(ru\leftarrow\sum_{z\in\tail_T(Y)}w(z)\)
    \STATE Append \(\langle tid,last,eu,ru\rangle\) to \(UL(Y)\)
\ENDFOR
\RETURN \(UL(Y)\)
\end{algorithmic}
\end{algorithm}

% =========================================================
\section{Correctness Analysis}
% =========================================================

\begin{theorem}[Soundness]
Every pattern emitted by HIEP-Miner belongs to \(\HIEP_{\theta,\sigma_0}\).
\end{theorem}

\begin{proof}
A pattern \(X\) is emitted only after HIEP-Miner computes its exact information utility:
\[
U(X)=\sum_{e\in UL(X)}e.eu.
\]
The algorithm emits \(X\) only if:
\[
U(X)\ge\theta.
\]
Furthermore, every emitted pattern has already passed:
\[
|UL(X)|=\supp(X)\ge\sigma_0.
\]
Therefore, every emitted pattern satisfies the HIEPM definition.
\end{proof}

\begin{theorem}[Completeness]
HIEP-Miner emits every pattern in \(\HIEP_{\theta,\sigma_0}\).
\end{theorem}

\begin{proof}
The depth-first search enumerates every pattern consistent with the fixed global item order unless a pruning rule stops its branch. Support pruning is safe by support anti-monotonicity. TIUB pruning is safe by the TIUB theorem. IWRU pruning is safe by the IWRU theorem. Therefore, if a pattern \(Y\) satisfies \(U(Y)\ge\theta\) and \(\supp(Y)\ge\sigma_0\), no prefix of \(Y\) can be pruned by these rules. Hence HIEP-Miner reaches and emits \(Y\).
\end{proof}

\begin{theorem}[Exactness of Utility-List Computation]
For every generated pattern \(X\), \(UL(X)\) contains exactly the supporting transactions of \(X\), and the sum of exact utility entries equals \(U(X)\).
\end{theorem}

\begin{proof}
The singleton lists are exact by construction. For an extension \(Y=X\cup\{a\}\), \textsc{Join-HIEP} retains exactly transactions that support both \(X\) and \(a\), and in sequence mode it additionally requires the correct order constraint. Thus the resulting list contains exactly \(\Gamma(Y)\). For each retained transaction, the algorithm computes \(eu=u(Y,T)\). Therefore:
\[
\sum_{e\in UL(Y)}e.eu=U(Y).
\]
\end{proof}

% =========================================================
\section{Complexity Analysis}
% =========================================================

Let \(n=|\D|\) be the number of window transactions, \(m=|\I|\) be the number of distinct tokens, and:
\[
nnz=\sum_{T\in\D}|T|
\]
be the number of token occurrences in the transaction database.

Let \(\Visited\) be the set of visited search nodes. For each node \(X\), let:
\[
s_X=|UL(X)|.
\]

\subsection{Preprocessing Complexity}

Computing token frequencies, probabilities, and weights requires:
\[
O(nnz).
\]

Building singleton vertical postings also requires:
\[
O(nnz).
\]

Computing transaction information utilities requires:
\[
O(nnz).
\]

Thus preprocessing time is:
\[
O(nnz).
\]

\subsection{Mining Complexity}

Each utility-list join is linear in the list sizes being merged:
\[
O(s_X+s_a).
\]

The total mining time is:
\[
O\left(
\sum_{X\in\Visited}
\sum_{Y\in children(X)}
JOIN(X,Y)
\right).
\]

In the worst case, the number of visited nodes is exponential in \(m\), because exact pattern mining is worst-case exponential. However, practical runtime is output-sensitive and pruning-sensitive: it depends on the number of prefixes surviving support, TIUB, and IWRU pruning.

\subsection{Memory Complexity}

The transaction store requires:
\[
O(nnz+n).
\]

The singleton postings require:
\[
O(nnz).
\]

The utility-list arena requires:
\[
O\left(
\max_{\pi}
\sum_{X\in\pi}s_X
\right),
\]
where \(\pi\) is an active depth-first branch.

Thus total memory is:
\[
O\left(
nnz+n+
\max_{\pi}\sum_{X\in\pi}s_X
+
|\Output|
\right).
\]

Since the utility-list arena is reused depth-first, HIEP-Miner avoids per-node heap allocation.

% =========================================================
\section{Experimental Design}
% =========================================================

\subsection{Research Questions}

A Q1-level experimental evaluation should address the following research questions:

\begin{itemize}
    \item RQ1: Does HIEP-Miner suppress high-frequency linguistic noise better than support-based miners?
    \item RQ2: Does HIEP-Miner recover planted rare high-information patterns with high recall?
    \item RQ3: How much search-space reduction is obtained by TIUB and IWRU?
    \item RQ4: How does HIEP-Miner compare with Apriori, FP-Growth, PrefixSpan, HUI-Miner, and EFIM in runtime and memory?
    \item RQ5: Can the embedded \texttt{mmap} tokenizer and miner pipeline process massive text streams at high throughput?
\end{itemize}

\subsection{Datasets}

The evaluation should use both real and synthetic datasets.

\begin{itemize}
    \item \textbf{Zipfian synthetic corpus}: generated text with stopword noise, multilingual tokens, and planted rare semantic patterns.
    \item \textbf{Retail}: sparse transaction benchmark.
    \item \textbf{Accidents}: dense transaction benchmark.
    \item \textbf{Chess}: dense structured benchmark.
    \item \textbf{Text logs}: command logs, chat logs, or software issue traces converted into token windows.
\end{itemize}

\subsection{Synthetic Corpus Model}

Let \(V\) be the vocabulary size. Token rank \(r\) follows the Zipf-Mandelbrot distribution:
\[
P(r)=
\frac{1}{Z}
\frac{1}{(r+q)^s},
\]
where:
\[
Z=\sum_{r=1}^{V}\frac{1}{(r+q)^s},
\]
\(s>1\) controls skew, and \(q\ge0\) is a shift parameter.

The synthetic corpus should contain:

\begin{enumerate}[label=(\arabic*)]
    \item high-frequency stopword filler;
    \item mid-frequency technical terms;
    \item low-frequency multilingual rare tokens;
    \item planted multi-token signal patterns;
    \item punctuation and casing noise.
\end{enumerate}

\subsection{Baselines}

\begin{table}[H]
\centering
\caption{Recommended baselines for evaluating HIEP-Miner.}
\label{tab:baselines}
\small
\begin{adjustbox}{max width=\textwidth}
\begin{tabular}{lll}
\toprule
Category & Baseline & Purpose \\
\midrule
Frequent itemset mining & Apriori & candidate-generation support baseline \\
Frequent itemset mining & FP-Growth & compressed-tree support baseline \\
Vertical itemset mining & Eclat & vertical support baseline \\
Sequential pattern mining & PrefixSpan & sequence-support baseline \\
High-utility mining & HUI-Miner & utility-list baseline \\
High-utility mining & EFIM & strong utility-mining baseline \\
Ablation & HIEP without TIUB & evaluates coarse bound contribution \\
Ablation & HIEP without IWRU & evaluates remaining-information bound \\
Ablation & HIEP with uniform weights & evaluates information weighting \\
Ablation & HIEP without compactness & evaluates phrase locality \\
\bottomrule
\end{tabular}
\end{adjustbox}
\end{table}

\subsection{Metrics}

\begin{definition}[Throughput]
\[
Throughput_{MB/s}
=
\frac{\text{bytes processed}/2^{20}}{\Delta t}.
\]
\end{definition}

\begin{definition}[Token Throughput]
\[
Throughput_{tok/s}
=
\frac{N}{\Delta t}.
\]
\end{definition}

\begin{definition}[Noise-Filtering Efficiency]
Let \(\Junk\) be the planted junk vocabulary and \(\Output\) be the output pattern set. Define:
\[
NFE=
1-
\frac{
|\{X\in\Output\mid X\subseteq\Junk\}|
}{
|\Output|
}.
\]
\end{definition}

\begin{definition}[Signal Recall]
Let \(\Signal\) be the set of planted signal patterns. Define:
\[
SignalRecall=
\frac{|\Output\cap\Signal|}{|\Signal|}.
\]
\end{definition}

\begin{definition}[Information Density Optimization]
\[
IDO=
\frac{1}{|\Output|}
\sum_{X\in\Output}
\frac{U(X)}{|X|}.
\]
\end{definition}

Additional system metrics include:

\begin{itemize}
    \item peak resident set size;
    \item number of visited nodes;
    \item number of generated candidates;
    \item number of branches pruned by support;
    \item number of branches pruned by TIUB;
    \item number of branches pruned by IWRU;
    \item output pattern count;
    \item output disk size;
    \item average pattern length;
    \item average pattern support;
    \item average information utility.
\end{itemize}

\subsection{Measurement Protocol}

Runtime should be measured with a monotonic clock. Peak memory should be measured using process-level peak resident memory. Each configuration should be repeated at least 11 times. Report the median and interquartile range.

For memory-mapped file scans, cold-cache and warm-cache experiments should be reported separately. Sequential-read hints should be applied in systems where they are available.

\subsection{Ablation Studies}

The following ablations are required:

\begin{enumerate}[label=(A\arabic*)]
    \item HIEP-Miner without compactness \(\kappa\).
    \item HIEP-Miner without TIUB.
    \item HIEP-Miner without IWRU.
    \item HIEP-Miner with uniform token weights.
    \item HIEP-Miner with different persistence floors \(\sigma_0\).
    \item HIEP-Miner with different Zipf skew parameters \(s\).
\end{enumerate}

\subsection{Expected Visualizations}

The paper should include:

\begin{itemize}
    \item runtime versus threshold \(\theta\);
    \item peak memory versus threshold \(\theta\);
    \item output pattern count versus \(\theta\);
    \item pruning breakdown stacked bar chart;
    \item Noise-Filtering Efficiency versus \(\theta\);
    \item Signal Recall versus \(\theta\);
    \item Information Density Optimization versus \(\theta\);
    \item throughput in MB/s and tokens/s;
    \item average and maximum utility-list length;
    \item effect of Zipf skew on output quality.
\end{itemize}

% =========================================================
\section{Experimental Evaluation}
% =========================================================

\subsection{Implementation and Reproducibility}

The results in this section were generated by the C implementation integrated into the unified framework executable:

\begin{verbatim}
bin/dm.exe hiep --input <dataset> --input-type text|transactions
              --mode itemset|sequence --minsup <ratio|count>
              --theta-ratio <value> --window <L> --stride <Delta>
              --tokenizer faro
\end{verbatim}

The experiment driver is \texttt{scripts/run\_hiep\_experiments.py}. It runs HIEP-Miner through \texttt{dm.exe hiep} on three transaction benchmarks, one real text corpus, four planted Zipfian text streams, and five ablation settings. External support-mining baselines are run by \texttt{scripts/run\_hiep\_baseline\_experiments.py}. The raw logs and parsed summaries are stored in \texttt{results/hiep\_q1/}; the figures included below are generated by \texttt{scripts/plot\_hiep\_results.py} and the baseline runner. Unless otherwise stated, the tokenizer is the FARO tokenizer exposed through \texttt{include/tokenizer/tokenizer.h}.

\subsection{Workloads}

\begin{table}[H]
\centering
\caption{Datasets and workload sizes used in the executed HIEP-Miner evaluation. Text corpora are converted to window transactions; transaction benchmarks are read as integer itemsets.}
\label{tab:hiep_workloads}
\small
\begin{adjustbox}{max width=\textwidth}
\begin{tabular}{llrrrrr}
\toprule
Dataset & Mode & Transactions & Stream tokens / nnz & Vocabulary & Window & Minimum support \\
\midrule
Retail & itemset & 12,000 & 120,896 & 9,004 & -- & 0.015 \\
Accidents & itemset & 3,500 & 118,457 & 283 & -- & 0.030 \\
Chess & itemset & 3,196 & 118,252 & 75 & -- & 0.080 \\
Databricks Dolly 15k & itemset text & 24,406 & 780,964 & 43,669 & 64/32 & 0.010 \\
Zipf \(s=1.05\) & sequence text & 3,750 & 90,000 & 670 & 48/24 & 3 \\
Zipf \(s=1.15\) & sequence text & 3,750 & 90,000 & 670 & 48/24 & 3 \\
Zipf \(s=1.25\) & sequence text & 3,750 & 90,000 & 670 & 48/24 & 3 \\
Zipf \(s=1.35\) & sequence text & 3,750 & 90,000 & 670 & 48/24 & 3 \\
\bottomrule
\end{tabular}
\end{adjustbox}
\end{table}

\subsection{Real Dataset Results}

Table~\ref{tab:hiep_real_results} reports representative real-dataset results at \(\theta\)-ratio \(0.50\), using the same depth and resource limits as the experiment scripts. HIEP-Miner runs within \(0.051\) seconds on Retail, \(1.402\) seconds on Accidents, \(2.018\) seconds on Chess, and \(7.938\) seconds on the 5 MB Dolly text slice. All four runs finish without hitting the time limit. The text run emits 3,290 high-information patterns from a 43,669-token vocabulary while maintaining \(41.07\) MB peak memory.

\begin{table}[H]
\centering
\caption{Representative HIEP-Miner results on real datasets at \(\theta\)-ratio \(0.50\).}
\label{tab:hiep_real_results}
\small
\begin{adjustbox}{max width=\textwidth}
\begin{tabular}{lrrrrrrr}
\toprule
Dataset & Runtime (s) & Peak RAM (MB) & Output & Visited nodes & IWRU pruned & IDO & NFE \\
\midrule
Retail & 0.051 & 6.590 & 14 & 73 & 50 & 4,689.78 & 1.000 \\
Accidents & 1.402 & 8.824 & 22,147 & 41,632 & 9,398 & 521.13 & 1.000 \\
Chess & 2.018 & 9.180 & 24,518 & 25,136 & 236 & 514.30 & 1.000 \\
Databricks Dolly 15k & 7.938 & 41.070 & 3,290 & 26,846 & 19,332 & 5,769.85 & 1.000 \\
\bottomrule
\end{tabular}
\end{adjustbox}
\end{table}

The Dolly threshold sweep shows the intended precision--selectivity behavior. As the \(\theta\)-ratio increases from \(0.20\) to \(0.70\), output count drops from 23,842 to 1,522 patterns, runtime drops from 14.484 seconds to 6.774 seconds, and information density rises from 2,567.58 to 7,644.28. This is the desired operating mode for paper-scale text mining: increasing the threshold compresses the output while retaining high-information patterns.

\subsection{External Baseline Comparison}

To avoid evaluating HIEP-Miner only against its own ablations, we compare against three support-based baselines already implemented in \texttt{dm.exe}: Apriori, Eclat, and FP-Growth. Each baseline is executed through the same framework command path on the same bounded transaction slices used by HIEP-Miner. The timeout is set to 35 seconds, matching the HIEP experiment limit.

This comparison is intentionally conservative and should be interpreted as a scalability and output-behavior comparison rather than an identical-objective comparison: Apriori, Eclat, and FP-Growth mine all support-frequent itemsets, while HIEP-Miner mines high-information patterns under \(\theta\), \(\sigma_0\), and depth constraints. Nevertheless, these algorithms are the standard external baselines for support-driven pattern mining, and they show the practical cost of relying on support alone in dense data.

\begin{table}[H]
\centering
\caption{External baseline comparison through \texttt{dm.exe}. HIEP uses \(\theta\)-ratio \(0.50\) and max depth 3. Support baselines use the same transaction slice and minimum support. TO denotes timeout at 35 seconds.}
\label{tab:hiep_external_baselines}
\small
\begin{adjustbox}{max width=\textwidth}
\begin{tabular}{llrrrl}
\toprule
Dataset & Algorithm & Runtime (s) & Peak RAM (MB) & Output count & Status \\
\midrule
Retail & HIEP-Miner & 0.051 & 6.590 & 14 & OK \\
Retail & Apriori & 0.040 & 4.450 & 115 & OK \\
Retail & Eclat & 0.013 & 5.680 & 97 & OK \\
Retail & FP-Growth & 0.005 & 4.850 & 97 & OK \\
\midrule
Accidents & HIEP-Miner & 1.402 & 8.824 & 22,147 & OK \\
Accidents & Apriori & \(>35.000\) & -- & -- & TO \\
Accidents & Eclat & \(>35.000\) & -- & -- & TO \\
Accidents & FP-Growth & \(>35.000\) & -- & -- & TO \\
\midrule
Chess & HIEP-Miner & 2.018 & 9.180 & 24,518 & OK \\
Chess & Apriori & \(>35.000\) & -- & -- & TO \\
Chess & Eclat & \(>35.000\) & -- & -- & TO \\
Chess & FP-Growth & \(>35.000\) & -- & -- & TO \\
\bottomrule
\end{tabular}
\end{adjustbox}
\end{table}

The sparse Retail slice is easy for classical support miners, and FP-Growth is fastest. On dense Accidents and Chess slices, however, all three support baselines fail to complete within the 35-second budget, while HIEP-Miner completes in 1.402 seconds and 2.018 seconds, respectively. The difference is explained by the HIEP objective and pruning design: support-only miners must enumerate a large dense frequent-pattern space, whereas HIEP-Miner uses self-information, TIUB, IWRU, and depth-bounded utility-list search to focus on high-information outputs.

\subsection{Planted Zipfian Signal Recovery}

Table~\ref{tab:hiep_zipf_results} reports sequence-mode results on planted Zipfian corpora at \(\theta\)-ratio \(0.35\). HIEP-Miner recovers all planted rare signal families for \(s=1.05\) and \(s=1.15\), and recovers \(75\%\) of them for the heavier-skew \(s=1.25\) and \(s=1.35\) corpora. Runtime decreases as skew increases because the search space contracts from 67,651 visited nodes at \(s=1.05\) to 21,119 visited nodes at \(s=1.35\).

\begin{table}[H]
\centering
\caption{Planted Zipfian sequence results at \(\theta\)-ratio \(0.35\).}
\label{tab:hiep_zipf_results}
\small
\begin{adjustbox}{max width=\textwidth}
\begin{tabular}{lrrrrrrr}
\toprule
Corpus & Runtime (s) & Peak RAM (MB) & Output & Visited nodes & IDO & NFE & Signal recall \\
\midrule
Zipf \(s=1.05\) & 1.134 & 10.281 & 251 & 67,651 & 927.19 & 0.542 & 1.000 \\
Zipf \(s=1.15\) & 1.159 & 9.715 & 249 & 49,662 & 849.41 & 0.434 & 1.000 \\
Zipf \(s=1.25\) & 0.891 & 9.227 & 245 & 33,678 & 795.98 & 0.343 & 0.750 \\
Zipf \(s=1.35\) & 0.527 & 8.488 & 223 & 21,119 & 771.09 & 0.269 & 0.750 \\
\bottomrule
\end{tabular}
\end{adjustbox}
\end{table}

\subsection{Ablation Results}

The ablation study in Table~\ref{tab:hiep_ablation_results} uses the Zipf \(s=1.15\) sequence corpus with \(\theta\)-ratio \(0.35\). Removing IWRU increases runtime from \(1.005\) seconds to \(6.532\) seconds, increases visited nodes by \(5.48\times\), and increases generated children by \(10.94\times\). This confirms that IWRU is the dominant exact branch-and-bound pruning rule. Removing TIUB has little effect in this workload because singleton support and IWRU already eliminate most branches. Uniform weights lose all planted signal recall, showing that Shannon self-information is essential rather than cosmetic. Removing compactness keeps recall but emits \(4.91\times\) more patterns and lowers NFE, confirming that locality control suppresses diffuse sequence matches.

\begin{table}[H]
\centering
\caption{Ablation results on Zipf \(s=1.15\), sequence mode, \(\theta\)-ratio \(0.35\).}
\label{tab:hiep_ablation_results}
\small
\begin{adjustbox}{max width=\textwidth}
\begin{tabular}{lrrrrrrr}
\toprule
Variant & Runtime (s) & Visited nodes & Generated children & IWRU pruned & Output & NFE & Signal recall \\
\midrule
Full HIEP-Miner & 1.005 & 49,662 & 230,467 & 47,496 & 249 & 0.434 & 1.000 \\
No TIUB & 1.145 & 49,662 & 230,467 & 47,496 & 249 & 0.434 & 1.000 \\
No IWRU & 6.532 & 272,185 & 2,520,460 & 0 & 249 & 0.434 & 1.000 \\
Uniform weights & 0.509 & 17,358 & 21,037 & 16,506 & 90 & 0.000 & 0.000 \\
No compactness & 0.845 & 49,662 & 230,467 & 47,496 & 1,223 & 0.293 & 1.000 \\
\bottomrule
\end{tabular}
\end{adjustbox}
\end{table}

\subsection{Generated Q1 Figures}

Figures~\ref{fig:hiep_external_baseline_runtime}, \ref{fig:hiep_system_charts}, and~\ref{fig:hiep_quality_charts} include the generated chart artifacts used for the experimental report. The source SVG files are under \texttt{results/hiep\_q1/}. When compiling with \texttt{pdflatex}, use \texttt{--shell-escape} for the \texttt{svg} package, or convert the SVG files to PDF/PNG and replace \texttt{\textbackslash includesvg} with \texttt{\textbackslash includegraphics}.

\begin{figure}[H]
\centering
\includesvg[width=0.92\textwidth,inkscapelatex=false]{../../results/hiep_q1/hiep_external_baseline_runtime}
\caption{Runtime comparison between HIEP-Miner and external support-mining baselines. Bars marked TO reached the 35-second timeout.}
\label{fig:hiep_external_baseline_runtime}
\end{figure}

\begin{figure}[H]
\centering
\begin{subfigure}{0.48\textwidth}
    \centering
    \includesvg[width=\linewidth,inkscapelatex=false]{../../results/hiep_q1/hiep_throughput}
    \caption{Tokenizer and miner throughput.}
\end{subfigure}
\hfill
\begin{subfigure}{0.48\textwidth}
    \centering
    \includesvg[width=\linewidth,inkscapelatex=false]{../../results/hiep_q1/hiep_pruning_breakdown}
    \caption{Support, TIUB, and IWRU pruning.}
\end{subfigure}
\caption{System-level HIEP-Miner performance charts generated from the executed experiments.}
\label{fig:hiep_system_charts}
\end{figure}

\begin{figure}[H]
\centering
\begin{subfigure}{0.48\textwidth}
    \centering
    \includesvg[width=\linewidth,inkscapelatex=false]{../../results/hiep_q1/hiep_zipf_skew_quality}
    \caption{Zipf skew and output quality.}
\end{subfigure}
\hfill
\begin{subfigure}{0.48\textwidth}
    \centering
    \includesvg[width=\linewidth,inkscapelatex=false]{../../results/hiep_q1/hiep_ablation_study}
    \caption{Ablation study.}
\end{subfigure}
\caption{Quality and ablation charts for the planted-signal evaluation.}
\label{fig:hiep_quality_charts}
\end{figure}

Overall, the executed results support the core claims of HIEPM: self-information weighting is required for rare-signal recovery, compactness reduces diffuse sequence noise, IWRU provides the main practical pruning power needed for exact high-information pattern mining, and external support-mining baselines struggle on dense slices where high-information pruning remains tractable.

% =========================================================
\section{Discussion}
% =========================================================

HIEPM changes the objective of text pattern mining. Instead of assuming that frequent means meaningful, it explicitly assigns low value to ubiquitous tokens and high value to specific tokens. This makes it better aligned with text logs, chats, command histories, and technical corpora where the most frequent tokens are often operational noise.

The key modeling decision is to separate the information-theoretic interpretation from the search-time objective. Direct joint surprisal is attractive but not directly mineable with classical downward closure. HIEPM therefore uses an additive entropy-derived utility that preserves exact mining through TIUB and IWRU upper bounds. This mirrors the historical development of HUIM, where utility mining became practical only after safe upper bounds were introduced.

The embedded systems design is equally important. HIEP-Miner is intended to run after a zero-dependency tokenizer that maps raw strings to integer IDs. The mining core does not need string comparisons, dynamic object allocation, or pointer-heavy structures. This makes it suitable for low-level text analytics engines such as \texttt{dm.exe}.

% =========================================================
\section{Threats to Validity}
% =========================================================

The first threat is corpus-relative weighting. Self-information depends on the empirical token distribution, so weights may drift in streaming settings. Incremental HIEPM with changing weights is a future research problem.

The second threat is rare-event bias. Pure information-theoretic scores can overvalue one-off rare artifacts. HIEPM addresses this with a persistence threshold \(\sigma_0\), compactness \(\kappa\), and planted-pattern validation, but further semantic filtering may be useful.

The third threat is tokenizer dependence. Different tokenization rules can produce different token universes and therefore different information weights. The tokenizer policy must be fixed and reported.

The fourth threat is baseline fairness. Apriori, FP-Growth, PrefixSpan, HUI-Miner, and EFIM solve related but not identical tasks. Comparisons should therefore report both runtime and output-quality metrics.

% =========================================================
\section{Conclusion}
% =========================================================

This paper introduced High-Information Entropy Pattern Mining, a new pattern mining paradigm for embedded text-stream analytics. HIEPM treats tokenized text as an integer transaction or sequence database and uses Shannon self-information as an entropy-derived utility signal. The resulting objective suppresses high-frequency linguistic noise and emphasizes rare, recurring, informative token combinations.

We formalized the HIEPM problem, proved that direct joint-surprisal objectives are non-monotone, and introduced a decomposable semantic information utility. We derived Transaction Information Upper Bound and Information-Weighted Remaining Utility as safe pruning bounds. We then proposed HIEP-Miner, a flat-array vertical utility-list algorithm designed for a zero-dependency C99-style engine.

The proposed framework bridges information retrieval, information theory, high-utility pattern mining, and high-performance embedded text analytics. Future work includes online HIEPM with drifting token distributions, closed high-information pattern mining, top-\(k\) HIEPM, privacy-preserving HIEPM, and contextual language-model-based reranking of exact HIEP outputs.

% =========================================================
% References
% =========================================================

\bibliographystyle{plain}
\begin{thebibliography}{99}

\bibitem{Shannon1948}
Claude E. Shannon.
\newblock A mathematical theory of communication.
\newblock \emph{Bell System Technical Journal}, 27(3):379--423, 1948.

\bibitem{Zipf1949}
George K. Zipf.
\newblock \emph{Human Behavior and the Principle of Least Effort}.
\newblock Addison-Wesley, 1949.

\bibitem{Piantadosi2014Zipf}
Steven T. Piantadosi.
\newblock Zipf's word frequency law in natural language: A critical review and future directions.
\newblock \emph{Psychonomic Bulletin \& Review}, 21:1112--1130, 2014.

\bibitem{AgrawalSrikant1994}
Rakesh Agrawal and Ramakrishnan Srikant.
\newblock Fast algorithms for mining association rules.
\newblock In \emph{Proceedings of the 20th International Conference on Very Large Data Bases}, pages 487--499, 1994.

\bibitem{HanPeiYin2000}
Jiawei Han, Jian Pei, and Yiwen Yin.
\newblock Mining frequent patterns without candidate generation.
\newblock In \emph{Proceedings of the ACM SIGMOD International Conference on Management of Data}, pages 1--12, 2000.

\bibitem{Zaki2000}
Mohammed J. Zaki.
\newblock Scalable algorithms for association mining.
\newblock \emph{IEEE Transactions on Knowledge and Data Engineering}, 12(3):372--390, 2000.

\bibitem{Pei2001PrefixSpan}
Jian Pei, Jiawei Han, Behzad Mortazavi-Asl, Helen Pinto, Qiming Chen, Umeshwar Dayal, and Mei-Chun Hsu.
\newblock PrefixSpan: Mining sequential patterns efficiently by prefix-projected pattern growth.
\newblock In \emph{Proceedings of the International Conference on Data Engineering}, pages 215--224, 2001.

\bibitem{ManningRaghavanSchutze2008IR}
Christopher D. Manning, Prabhakar Raghavan, and Hinrich Sch{\"u}tze.
\newblock \emph{Introduction to Information Retrieval}.
\newblock Cambridge University Press, 2008.

\bibitem{ChurchGale1995}
Kenneth W. Church and William A. Gale.
\newblock Poisson mixtures.
\newblock \emph{Natural Language Engineering}, 1(2):163--190, 1995.

\bibitem{TomokiyoHurst2003}
Takashi Tomokiyo and Matthew Hurst.
\newblock A language model approach to keyphrase extraction.
\newblock In \emph{Proceedings of the ACL Workshop on Multiword Expressions}, pages 33--40, 2003.

\bibitem{Tatti2008MaximumEntropy}
Nikolaj Tatti.
\newblock Maximum entropy based significance of itemsets.
\newblock \emph{Knowledge and Information Systems}, 17:57--77, 2008.

\bibitem{DeBie2011SubjectiveInterestingness}
Tijl De Bie.
\newblock An information theoretic framework for data mining.
\newblock In \emph{Proceedings of the ACM SIGKDD International Conference on Knowledge Discovery and Data Mining}, pages 564--572, 2011.

\bibitem{Kontonasios2012SubjectiveInterestingness}
Kleanthis-Nikolaos Kontonasios and Tijl De Bie.
\newblock Subjectively interesting pattern mining.
\newblock In \emph{Proceedings of the European Conference on Machine Learning and Principles and Practice of Knowledge Discovery in Databases}, pages 42--57, 2012.

\bibitem{Vreeken2011Krimp}
Jilles Vreeken, Matthijs van Leeuwen, and Arno Siebes.
\newblock Krimp: Mining itemsets that compress.
\newblock \emph{Data Mining and Knowledge Discovery}, 23:169--214, 2011.

\bibitem{Galbrun2020MDL}
Esther Galbrun.
\newblock The minimum description length principle for pattern mining: A survey.
\newblock \emph{Data Mining and Knowledge Discovery}, 34:1679--1727, 2020.

\bibitem{Liu2005TwoPhase}
Ying Liu, Wei-keng Liao, and Alok Choudhary.
\newblock A two-phase algorithm for fast discovery of high utility itemsets.
\newblock In \emph{Proceedings of the Pacific-Asia Conference on Knowledge Discovery and Data Mining}, pages 689--695, 2005.

\bibitem{Liu2012HUIMiner}
Mengchi Liu and Junfeng Qu.
\newblock Mining high utility itemsets without candidate generation.
\newblock In \emph{Proceedings of the ACM International Conference on Information and Knowledge Management}, pages 55--64, 2012.

\bibitem{FournierViger2014FHM}
Philippe Fournier-Viger, Cheng-Wei Wu, Souleymane Zida, and Vincent S. Tseng.
\newblock FHM: Faster high-utility itemset mining using estimated utility co-occurrence pruning.
\newblock In \emph{Proceedings of the International Symposium on Methodologies for Intelligent Systems}, pages 83--92, 2014.

\bibitem{Zida2017EFIM}
Souleymane Zida, Philippe Fournier-Viger, Jerry Chun-Wei Lin, Cheng-Wei Wu, and Vincent S. Tseng.
\newblock EFIM: A fast and memory efficient algorithm for high-utility itemset mining.
\newblock \emph{Knowledge and Information Systems}, 51:595--625, 2017.

\bibitem{Gan2021UtilitySurvey}
Wensheng Gan, Jerry Chun-Wei Lin, Philippe Fournier-Viger, Han-Chieh Chao, Vincent S. Tseng, and Philip S. Yu.
\newblock A survey of utility-oriented pattern mining.
\newblock \emph{IEEE Transactions on Knowledge and Data Engineering}, 33(4):1306--1327, 2021.

\bibitem{FournierViger2014SPMF}
Philippe Fournier-Viger, Antonio Gomariz, Ted Gueniche, Azadeh Soltani, Cheng-Wei Wu, and Vincent S. Tseng.
\newblock SPMF: A Java open-source pattern mining library.
\newblock \emph{Journal of Machine Learning Research}, 15:3389--3393, 2014.

\bibitem{FIMIRepository}
Bart Goethals.
\newblock Frequent Itemset Mining Implementations Repository.
\newblock 2003.

\bibitem{Yergeau2003UTF8}
Fran\c{c}ois Yergeau.
\newblock UTF-8, a transformation format of ISO 10646.
\newblock RFC 3629, Internet Engineering Task Force, 2003.

\bibitem{KerriskMmap}
Michael Kerrisk.
\newblock mmap(2): Linux manual page.
\newblock Linux man-pages project.

\bibitem{KerriskClock}
Michael Kerrisk.
\newblock clock\_gettime(3): Linux manual page.
\newblock Linux man-pages project.

\bibitem{KerriskGetrusage}
Michael Kerrisk.
\newblock getrusage(2): Linux manual page.
\newblock Linux man-pages project.

\bibitem{KerriskProc}
Michael Kerrisk.
\newblock proc\_pid\_status(5): Linux manual page.
\newblock Linux man-pages project.

\end{thebibliography}

\end{document}
